\documentclass[11pt]{article}
\usepackage[margin=0.9in]{geometry}
\usepackage{graphicx,amsmath,amssymb,booktabs,hyperref,xcolor}
\hypersetup{colorlinks=true,linkcolor=blue,urlcolor=blue}
\setlength{\parskip}{6pt}
\title{Pad\'e vs Richardson for the K=3 lookup function:\\
       rational extrapolation does NOT beat polynomial}
\author{Fixed-Point Factory --- two-kernel extrapolation test}
\date{June 2026}

\begin{document}
\maketitle

\begin{abstract}
\noindent
Tested: replace Richardson (polynomial in $h^2$) with Pad\'e rational
approximants Pad\'e(0,1) $=\frac{a}{1+b h^2}$ and Pad\'e(1,1)
$=\frac{a + b h^2}{1 + c h^2}$ to extrapolate the kernel-band lookup
to $h{=}0$. Hypothesis: rational forms handle non-analytic $h$-dependence
better than polynomial.

\textbf{Result: hypothesis wrong.} On 16 saved R4 FPs at $G{=}11$:

\begin{itemize}
  \item At $\tau{=}0.2$ (smooth regime): Pad\'e(0,1) ties Richardson R2
        ($\max|\Delta|=0.21$); Pad\'e(1,1) DEGENERATES because
        $\mu \approx 0.5$ over flat regions makes the rational
        denominator's linear system rank-deficient -- returns
        constant 0.5.
  \item At $\tau{=}1$ (kinky regime): Pad\'e(0,1) max error
        $\sim 0.80$ to $0.91$ (vs Richardson R2 max $\sim 0.4$);
        Pad\'e(1,1) max $\sim 0.67$. Both WORSE than Richardson.
\end{itemize}

\textbf{Mechanism.} The kernel-band $\mu(h)$ does \emph{not} have poles
in $h^2$. It has \emph{jumps} -- non-analytic discontinuities at $h$ values
where the band edges sweep across kinks/critical points. Pad\'e
approximants try to fit poles where there are none; the denominator
$1 + b h^2$ approaches zero spuriously, giving wild extrapolations.
Richardson's polynomial form is safer (can't blow up).

\textbf{Conclusion.} The two-kernel extrapolation idea is mechanically
sound but doesn't help because the underlying $\mu(h)$ isn't pole-like.
The right path to handle the non-analyticity is still
\textbf{cusp-aware quadrature} (handle each topology change as an
analytic singular contribution), not a rational approximant in $h$.
\end{abstract}

\section{Cell-by-cell results}

\begin{center}
\begin{tabular}{lrrrrr}
\toprule
cell & Rich R2 max & Rich R4 max & Pad\'e(0,1) max & Pad\'e(1,1) max \\
\midrule
$\gamma{=}10$, $\tau{=}0.2$  & 0.21 & 0.21 & 0.21 & 0.50 \\
$\gamma{=}10$, $\tau{=}1$    & 0.43 & 0.57 & 0.91 & 0.68 \\
$\gamma{=}100$, $\tau{=}0.2$ & 0.21 & 0.21 & 0.21 & 0.50 \\
$\gamma{=}100$, $\tau{=}1$   & 0.38 & 0.44 & 0.85 & 0.68 \\
$\gamma{=}1000$, $\tau{=}0.2$ & 0.21 & 0.21 & 0.21 & 0.50 \\
$\gamma{=}1000$, $\tau{=}1$  & 0.38 & 0.44 & 0.81 & 0.68 \\
\bottomrule
\end{tabular}
\end{center}

\section{Why Pad\'e fails}

For Pad\'e(0,1) $f(h) = a / (1 + b h^2)$ fit to two points:
\begin{align*}
b &= (A_2 - A_1) / (A_1 h_1^2 - A_2 h_2^2) \\
a &= A_1 (1 + b h_1^2)
\end{align*}
If $A_1, A_2$ are close (smooth $\mu$), the numerator $A_2 - A_1$ is
small AND the denominator $A_1 h_1^2 - A_2 h_2^2 \approx
A_\mathrm{common}(h_1^2 - h_2^2)$ is the difference of two close
numbers --- catastrophic cancellation. The estimator is noisy.

If $\mu$ has true poles (which it doesn't here, only jumps),
$b > 0$ captures them and extrapolation is exact. But for jump
discontinuities, the data look like a low-degree polynomial PLUS noise
from the jump; Pad\'e treats the noise as a pole and extrapolates
into a spurious singularity.

For Pad\'e(1,1) with 3 points: 3$\times$3 linear system in $(a, b, c)$
becomes near-singular when the three $A_i$ are nearly equal (flat $\mu$);
\texttt{np.linalg.solve} returns garbage that we fall back to mean of
inputs, hence the 0.5 result at smooth cells.

\section{Conclusion}

Pad\'e in $h^2$ is the wrong tool for this operator. The right path
remains:
\begin{enumerate}
  \item Detect critical points of $P$ analytically (zeros of $\nabla P$,
        solvable from Cheby coefs via 2D resultant).
  \item Add an analytic kink residual to the GL quadrature near each
        critical $p^*$.
\end{enumerate}
This makes the operator continuous in $P$ and restores exponential
spectral convergence at any $\tau$.

\section{Reproducibility}
\texttt{projects/REZN/solver\_code/highprec\_dd\_qd/dd\_k3\_pade.py}

\end{document}
